\section{Chat benchmarks：开放式回答如何评价}

真实用户不会总问多选题。用户会问“甜菜和山羊奶酪做沙拉配什么香草”，这种问题没有唯一答案，甚至好答案依赖口味、上下文和表达风格。因此开放式评价要从 accuracy 转向 preference、rubric 和 judge reliability。

\begin{figure}[H]
\centering
\includegraphics[width=0.88\textwidth]{arena-beets.png}
\caption{Chatbot Arena 示例：开放式 beet salad prompt，两模型匿名回答，由用户选择更好者。}
\figsource{CS336 Lecture 12 官方源引用图。}
\end{figure}

\begin{knowledgebox}{读图：开放式评价为什么难}
同一个 prompt 可能有多个好答案。回答可以在正确性、风格、详细程度、简洁性、礼貌性上不同。单一 accuracy 无法覆盖这些维度，因此 pairwise preference 成为高信号替代。图中例子看似日常，却正好暴露“好回答”不是一个标量事实。
\end{knowledgebox}

\[
P(A \text{ beats } B)
  =
\frac{1}{1+10^{(ELO_B-ELO_A)/400}}.
\]
这里 \(ELO_A,ELO_B\) 是两个模型的评分。Arena 用大量 pairwise comparisons 拟合 ELO，从而得到动态 leaderboard。

\begin{figure}[H]
\centering
\includegraphics[width=0.58\textwidth]{lmarena-leaderboard.png}
\caption{LM Arena leaderboard：基于 pairwise human preference 的模型排序。}
\figsource{CS336 Lecture 12 官方源引用图。}
\end{figure}

\begin{warningbox}{老师强调：Arena 真实但不干净}
Arena 的 prompts 来自真实用户，这是它的价值；但用户是谁、是否懂任务、是否偏好长回答、是否被 sycophancy 影响，都会改变结果。Arena 很真实，但真实不等于无偏。它适合回答“真实用户更喜欢哪个系统”，不适合单独回答“哪个模型更正确”。
\end{warningbox}

LLM judge 的出现，是因为人类偏好昂贵且慢。AlpacaEval 和 WildBench 试图用模型裁判和 rubric/checklist 降低成本，同时用 Chatbot Arena 相关性校准它们是否接近人类偏好。

\begin{figure}[H]
\centering
\includegraphics[width=0.9\textwidth]{alpacaeval-correlation.png}
\caption{AlpacaEval 与 Chatbot Arena 的相关性：用 LLM judge 近似人类偏好。}
\figsource{AlpacaEval 官方图，本地化后供 CS336 Lecture 12 使用。}
\end{figure}

\begin{figure}[H]
\centering
\includegraphics[width=0.62\textwidth]{alpacaeval-leaderboard.png}
\caption{AlpacaEval leaderboard：以 baseline model 为参照，用 judge 评估 win rate。}
\figsource{CS336 Lecture 12 官方源引用图。}
\end{figure}

\begin{knowledgebox}{读图：LLM judge 的优势和危险}
LLM judge 便宜、可扩展、可复现，但会偏好长回答、特定风格或自己熟悉的模型。AlpacaEval 2.0 用 regression debias，是因为指标本身也需要被评价。读相关性图时要注意：高相关不等于无偏，只说明在某个分布上它能近似 Arena。
\end{knowledgebox}

\begin{figure}[H]
\centering
\includegraphics[width=0.88\textwidth]{wildbench.png}
\caption{WildBench：从真实人机对话中抽样，用 checklist 和 GPT-4 judge 提高开放式评价可靠性。}
\figsource{CS336 Lecture 12 官方源引用图。}
\end{figure}

\begin{importantbox}{开放式评价的经验规则}
Pairwise comparison 通常比单独打分更稳定；checklist/rubric 能提高人类或 LLM judge 的一致性；必须监控 judge bias、长度偏好、style preference 和 correctness/style conflation。老师在这里给出的实践判断是：开放式评价不是不能做，而是必须显式管理偏差。
\end{importantbox}

\subsection{本章小结}

Chat benchmarks 更接近真实助理使用，但牺牲了可控性。它们测的是用户或 judge 在特定分布上的偏好，不是纯粹能力。下一节继续从“说得好”推进到“做得成”：agentic benchmarks。

